\chapter{Result and Evaluation}
\section{Result}
This chapter presents the experimental results obtained from deploying and testing the
RedSentinel system against a live intentionally vulnerable web application hosted on
PortSwigger Web Security Academy. The evaluation demonstrates the system's end-toend capability: from initiating a scan through the Next.js dashboard, to automated
vulnerability discovery, AI-driven context classification, payload injection and browser
verification, and finally structured report generation.
\subsection{Experimental Setup}
RedSentinel was evaluated against a real, publicly accessible vulnerable lab instance provided by PortSwigger Web Security Academy. The target URL was:
\smallskip
\begin{Verbatim}[fontsize=\footnotesize]
  https://0a9e00260413f8418289ac3d00e600d4.web-security-academy.net  
\end{Verbatim}
\smallskip
The scan was configured in Single Page mode with a maximum of 50 payloads per parameter, WAF bypass enabled, and verify execution set to true. The complete Docker Compose stack was deployed locally, comprising the NestJS core, three Python AI microservices (Context, Payload-Gen, Fuzzer), Redis, PostgreSQL, and the Next.js dashboard.

\subsection{Dashboard and Scan Initiation}
Figure A.1 shows the RedSentinel dashboard at the time of scan initiation. The dashboard
displays aggregate statistics across all historical scans: 15 completed scans and 4 total
vulnerabilities discovered. A new scan is configured against the target URL with scan
mode set to Single Page and the max payloads slider set to 50.


The interface provides real-time visibility into scan progress via WebSocket updates
(Socket.io). Once the Start Scan button is clicked, the Core API persists the scan record
to PostgreSQL, enqueues a BullMQ job, and the dashboard immediately begins receiving
scan:progress events as each pipeline phase completes.

\subsection{Scan Execution and Vulnerability Discovery}
Figure A.2 shows the Scan Detail view upon completion. The scan completed in 2 minutes
and 44 seconds, processing the target through all five pipeline phases: crawling, context
classification, payload generation, fuzzing, and report generation. A total of 2
vulnerabilities were identified.


The findings panel reports:
\begin{itemize}
    \item 1 CRITICAL severity finding — Stored Cross-Site Scripting (XSS) in the searchparameter, confirmed via headless browser execution (EXEC confirmed).
    \item 1 HIGH severity finding — SVG/Polyglot XSS in the comment parameter, detected via reflection analysis but classified as unconfirmed (suspected).
\end{itemize}

 
 
The scan configuration is also shown at the bottom of the detail view: depth = 1, maxParams = 100, verifyExecution = true, wafBypass = true, maxPayloadsPerParam = 50, timeout = 60000ms, reportFormat = html, json, pdf.

\subsection{Detailed Vulnerability Findings}
\subsubsection{ Issue \#1 — Stored Cross-Site Scripting (XSS) [CRITICAL]}
Figure A.3 presents the detailed breakdown for Issue \#1. The vulnerability was discovered
in the search input field on the application homepage. The DistilBERT context classifier
identified the injection context as HTML\_BODY, and the Payload-Gen module selected
and ranked an SVG-based payload. The Fuzzer module confirmed successful JavaScript
execution in Playwright headless Chromium, classifying this as a CONFIRMED,
CRITICAL severity finding.

\begin{table}[h]
\centering
\begin{tabular}{|l|p{10cm}|}
\hline
\textbf{Page} & https://0a9e00260413f8418289ac3d00e600d4.web-security-academy.net/ \\ \hline
\textbf{Input field} & search \\ \hline
\textbf{Risk level} & CRITICAL \\ \hline
\textbf{Test code} & \texttt{<svg onload=alert()>} \\ \hline
\textbf{Confirmed} & Yes — JavaScript execution verified in headless browser \\ \hline
\end{tabular}
\end{table}

The remediation recommendation generated by the report engine advises sanitizing all
user-submitted content before storage, encoding stored content on render, and
implementing a Content Security Policy (CSP) header.

\subsubsection{ Issue \#2 — SVG/Polyglot XSS [HIGH]}
Figure A.4 presents the detailed breakdown for Issue \#2. SVG content rendered with usercontrolled data in the comment field was found to allow injection of SVG namespace
events (such as onload, onerror) or nested script tags. SVG parsers operate within a
security context separate from standard HTML parsing, allowing payloads that bypass
HTML-only sanitization.

\begin{table}[h]
\centering
\begin{tabular}{|l|p{10cm}|}
\hline
\textbf{Page} & https://0a9e00260413f8418289ac3d00e600d4.web-securityacademy.net/post?postId=2 \\ \hline
\textbf{Input field} & comment \\ \hline
\textbf{Risk level} & HIGH \\ \hline
\textbf{Test code} & \texttt{alert(1);} \\ \hline
\textbf{Confirmed} & Not confirmed (suspected — reflection detected, browser execution not triggered) \\ \hline
\end{tabular}
\end{table}
The report engine's remediation advice for this finding includes: parsing and validating user-provided SVG strictly against an allowlist of elements and attributes; serving SVG files with Content-Type application/svg+xml rather than text/html; never rendering unsanitized SVG data via innerHTML or in an HTML context; and sanitizing SVG using a library such as svg-sanitizer or DOMPurify.

\subsection{Summary of Results}
Table 5.1 summarizes the complete findings from the evaluation scan.

\begin{table}[h]
\centering
\small
\setlength{\tabcolsep}{4pt}
\renewcommand{\arraystretch}{1.2}

\resizebox{\textwidth}{!}{
\begin{tabular}{|c|p{4.5cm}|c|c|c|c|}
\hline
\textbf{\#} &
\textbf{URL / Parameter} &
\textbf{Input Field} &
\textbf{Severity} &
\textbf{Type} &
\textbf{Confirmed} \\
\hline

1 &
\texttt{...academy.net/ (search)} &
search &
CRITICAL &
Stored XSS &
Yes \\
\hline

2 &
\texttt{...academy.net/post\allowbreak ?postId=2 (comment)} &
comment &
HIGH &
SVG/Polyglot XSS &
No \\
\hline

\end{tabular}
}

\caption{Summary of Vulnerabilities Discovered}

\end{table}

The system successfully completed the full end-to-end pipeline, from scan initiation
through vulnerability discovery, browser-verified execution confirmation, and multi-format
report generation, in under three minutes without any manual
intervention. The CRITICAL stored XSS finding was fully confirmed via headless
Chromium execution using Playwright, validating the dual-verification approach described
in the Methodology chapter. The HIGH-severity SVG/Polyglot XSS finding was surfaced
through reflection analysis and context classification but was not browser-confirmed in
this scan run, consistent with the system's classification of such findings as SUSPECTED
rather than CONFIRMED.

These results demonstrate that RedSentinel's AI-driven context classification and
payload generation pipeline is capable of discovering both common and context-specific
XSS vulnerabilities that traditional static-payload scanners may miss, while its headless
browser verification stage ensures that confirmed findings represent genuine security
risks rather than benign reflections.

\section{Evaluation}
